\chapter{\uppercase{Conclusion}}

This research establishes a foundation for an adaptive and intelligent multi-stage vulnerability detection framework that enhances the reliability and contextual awareness of Static Application Security Testing (SAST) systems. By addressing the persistent challenges of high false positives, limited scalability, and lack of confidence-aware triage, the proposed framework advances beyond conventional static analysis paradigms. The system’s modular and layered architecture allows efficient detection, prioritization, and interpretation of vulnerabilities, integrating seamlessly into CI/CD environments and DevSecOps workflows. The evaluation demonstrates that this approach achieves improved balance between detection precision and runtime performance, making it suitable for modern continuous-integration ecosystems where timely, reliable, and explainable feedback is essential.

Building upon these outcomes, \textbf{Phase 2} of this research introduces advanced reasoning and adaptive intelligence components designed to extend the system’s analytical depth and decision accuracy. The first enhancement involves the \textbf{Agentic Graph Neural Network (GNN)} model trained on \textbf{Code Property Graph (CPG)} representations, enabling structural reasoning and uncertainty estimation. This step allows the framework to understand deep relational patterns within source code, learn semantic dependencies, and quantify prediction confidence critical for prioritizing findings based on contextual risk. Through iterative training and reinforcement on curated vulnerability corpora, the GNN phase enhances both interpretability and scalability across diverse codebases.

The second enhancement focuses on \textbf{LLM integration with Retrieval-Augmented Generation (RAG)} for semantic validation and contextual grounding. Here, a domain-adapted large language model (LLM) interprets ambiguous or borderline cases identified by the GNN, using external knowledge retrieval to verify potential vulnerabilities against known patterns, CVEs, and secure coding guidelines. This mechanism introduces semantic cross-verification, allowing the system to reason over both structural and natural-language evidence before final classification.

Finally, the framework evolves toward a \textbf{multimodal fusion and active escalation architecture}, implementing a hierarchical cascade   \textbf{SAST $\rightarrow$ GNN $\rightarrow$ LLM}   governed by uncertainty thresholds. When a detection’s confidence level falls below a predefined threshold, control escalates dynamically to higher reasoning layers, ensuring optimal trade-offs between speed, depth, and reliability. This active fusion approach allows the system to autonomously self-regulate analysis depth based on contextual uncertainty, minimizing false alarms while retaining comprehensive coverage.

In conclusion, these Phase 2 extensions transform the framework into a robust, agentic, and explainable security analysis ecosystem capable of structural reasoning, semantic grounding, and adaptive escalation. Future work will emphasize large-scale evaluation of the multimodal cascade across polyglot codebases, the integration of continual learning for evolving threat landscapes, and the incorporation of explainable AI (XAI) modules for developer interpretability. The long-term vision is to realize a fully autonomous, self-learning security intelligence agent that continuously improves its detection fidelity and contextual awareness across the entire software development lifecycle.

